\section{The construction}
\label{sec:construction}

We construct a map with a circular clock, a main sphere coordinate, and
finitely many auxiliary sphere coordinates, called registers.  The main
coordinate usually squares.  Early in each clock period, the map inserts
small signals into the registers; later in that period, it adds the
amplified signals to the main coordinate.  The inserted signals predict
the phase of a future square.  Their coefficients decrease with the
predicted delay, which makes the resulting infinite series continuously
differentiable.  This section specifies the map and proves its regularity;
Section~\ref{sec:orbits} analyzes its orbits.

Put $\Sigma=\mathbb C\cup\{\infty\}$, with the smooth structure of the
Riemann sphere, and let $t$ denote a coordinate on
$\mathbb R/100\mathbb Z$.  We also use $t$ for a real lift when specifying
functions on one clock period.  All occurrences of $z\mapsto z^2$ on
$\Sigma$ include $\infty\mapsto\infty$.

\subsection{The clock and its intervals}

The clock moves by at most one at each iterate and can slow down only
near $t=50$.  Choose a smooth $100$-periodic function
$b\colon\mathbb R\to[0,1]$ such that
\begin{equation}
\label{eq:clock-barrier}
\begin{gathered}
 b(t)=0\quad\Longleftrightarrow\quad t=50\pmod{100},
 \qquad b(t)=1\quad\text{if }|t-50|\ge2\text{ in }[0,100],\\
 t+b(t)\le50\qquad(0\le t\le50).
\end{gathered}
\end{equation}
These requirements are compatible.  Indeed, let
$\eta\colon\mathbb R\to[0,1]$ be smooth, equal to one on $[-1,1]$,
and supported in $(-2,2)$.  On $[0,100]$ take
\[
 b(t)=\eta(t-50)\frac{(t-50)^2}{4}+1-\eta(t-50),
\]
and extend periodically.  For $x=50-t\in[0,1]$ the value is $x^2/4\le x$;
for $x\in[1,2]$ both $x^2/4$ and $1$ are at most $x$; and for $x\ge2$
the value is $1\le x$.  This proves the last condition in
\eqref{eq:clock-barrier}.  The other conditions follow directly from the
formula, which is constant near both endpoints of the period.

Choose a smooth $e\colon\Sigma\to[0,1]$ that is zero on
$\{|z|\le1/2\}$ and one on $\{|z|\ge3/4\}\cup\{\infty\}$, and set
\begin{equation}
\label{eq:clock-increment}
 h(t,z)=b(t)+(1-b(t))e(z).
\end{equation}
A radial cutoff in the affine chart, constant near zero and infinity,
gives such an $e$.  Thus $0\le h\le1$, $h\ge b$, and $h=1$ whenever
$b=1$ or $e=1$.  In particular, the increment is exactly one throughout
the preparation and readout intervals specified next.

Choose smooth functions with values in $[0,1]$ as follows:
\begin{enumerate}
\item $\alpha\colon\mathbb R\to[0,1]$ is supported in $(8,12)$
      and equals one on $[9,11]$.
\item For a finite integer $q\ge1$, the functions
      $\beta_\ell\colon\mathbb R\to[0,1]$, $1\le \ell\le q$, have supports
      contained in intervals $J_\ell\subset(74,80)$ of length less than one,
      and satisfy
      \begin{equation}
      \label{eq:readout-cover}
       \forall t\in[76,77]\quad
       \exists \ell\in\{1,\ldots,q\}\quad \beta_\ell(t)=1.
      \end{equation}
\item $\chi\colon\Sigma\to[0,1]$ is supported away from zero and infinity
      and equals one on $\{1/2\le |z|\le1\}$.
\end{enumerate}
For the second choice, cover $[76,77]$ by finitely many smaller intervals
whose closures lie in intervals $J_\ell$ as above, and choose $\beta_\ell$ equal
to one on those closures.  The first and third choices are ordinary
smooth cutoffs.  We write
\[
 \bar\beta_\ell(t)=\sum_{k\in\mathbb Z}\beta_\ell(t-100k)
\]
for the $100$-periodization.  The translated supports are disjoint, so
$0\le\bar\beta_\ell\le1$.  The functions
$\beta_\ell$ themselves, as opposed to $\bar\beta_\ell$, will be used for the
predictions within a single lifted period.
Figure~\ref{fig:clock} summarizes the order of these clock intervals
and the reset interval specified below.

\subsection{Prediction inputs and the update rule}

To predict a later clock reading, temporarily omit the additions to the
main coordinate and define the smooth map
\begin{equation}
\label{eq:prediction-map}
 G\colon\mathbb R\times\Sigma\longrightarrow\mathbb R\times\Sigma,
 \qquad G(t,z)=(t+h(t,z),z^2),
 \qquad T_m=\operatorname{pr}_1\circ G^m\quad(m\ge1).
\end{equation}
Fix a product Riemannian metric $dt^2+g_\Sigma$ on $\mathbb R\times\Sigma$.
The map $G$ commutes with translation by $100$ in the first coordinate.
Compactness of $[0,100]\times\Sigma$ therefore supplies a constant
$B\ge1$ with $\|DG\|\le B$ everywhere.  Fix
\begin{equation}
\label{eq:gain-choice}
 L>\max\{B,2\}.
\end{equation}
The two inequalities in \eqref{eq:gain-choice} will control the derivatives
of the clock prediction and of the angular power, respectively.

Each register must admit multiplication by $L$ near zero while remaining
a sphere coordinate globally.  Choose a smooth function
$p\colon[0,\infty)\to[0,3]$ with
\[
 p(r)=r\quad(0\le r\le1),\qquad
 p\text{ nondecreasing on }[0,3L],\qquad
 p(r)=0\quad\text{for all sufficiently large }r.
\]
For existence, integrate from $0$ to $r$ a smooth $[0,1]$-valued function
that equals one on $[0,1]$ and zero on $[2,\infty)$.  Its integral is nondecreasing,
agrees with $r$ up to $1$, and is constant with value at most $2$ beyond
$2$.  Multiplication by a smooth cutoff equal to one on $[0,3L]$ and
zero for sufficiently large radii gives the required $p$.  Define
\begin{equation}
\label{eq:register-saturation}
 S\colon\Sigma\longrightarrow\mathbb C,
 \qquad
 S(v)=p(|v|)\frac{v}{|v|}\quad(0<|v|<\infty),
 \qquad S(0)=S(\infty)=0.
\end{equation}
This map is smooth: it is the identity near zero and identically zero
near infinity.  Moreover, $|S(v)|\le3$ everywhere and $S(v)=v$ on the
closed unit disk.  Choose a smooth $100$-periodic function
$a\colon\mathbb R\to[0,L]$ such that
\begin{equation}
\label{eq:register-gain}
 a(t)=0\quad(t\in[2,5]),\qquad
 a(t)=L\quad(t\in[7,90]).
\end{equation}
For example, multiply $L$ by a smooth cutoff supported in $(6,96)$ and
equal to one on $[7,90]$, then extend periodically.

For $0\le t<100$ and $1\le \ell\le q$, define
\begin{equation}
\label{eq:input}
 U_\ell(t,z)=\frac{\alpha(t)}{10}
 \sum_{m=1}^{\infty} L^{1-m}\beta_\ell(T_m(t,z))\,
 \chi(z)\left(\frac{z}{|z|}\right)^{2^{m+1}}.
\end{equation}
The product containing $\chi$ is defined to be zero at $z=0,\infty$.
Extend $U_\ell$ itself periodically in $t$.  The series and its first
derivatives are proved to converge below.  Notice that the factor in the
series is the function $\beta_\ell$ on the real line, not its periodization:
only predicted readings in $J_\ell\subset(74,80)$ contribute.
The index $m$ counts iterates of $G$, not displacement of the clock.

Set
\[
 M=(\mathbb R/100\mathbb Z)\times\Sigma\times\Sigma^q,
\]
with coordinates $(t,z,v_1,\ldots,v_q)$.  Define $f\colon M\to M$ by
the simultaneous update
\begin{equation}
\label{eq:map}
\begin{aligned}
 t'&=t+h(t,z)\pmod{100},\\
 z'&=z^2+20\sum_{\ell=1}^q\bar\beta_\ell(t)S(v_\ell),\\
 v_\ell'&=a(t)S(v_\ell)+U_\ell(t,z),\qquad 1\le \ell\le q.
\end{aligned}
\end{equation}
We set $z'=\infty$ when $z=\infty$; every register output is in the
affine complex chart.  All right-hand sides use the old state.  In
particular, the values $t\in[2,5]$ reset every register to zero, because
both $a(t)$ and $U_\ell(t,z)$ vanish there.

Here is the timing behind \eqref{eq:input}.  Write $z_i$ for the main
coordinate at iteration $i$.  A signal evaluated at time $i$ enters
its register at time $i+1$.  If it is to be used in the main
update at time $i+m$, it receives $m-1$ intervening gains.  Provided
$a=L$ during those updates and the register remains in the unit disk,
the initial factor $L^{1-m}$ is exactly canceled by those gains.  If the
main coordinate follows squaring, its squared phase at time $i+m$ is
$(z_i/|z_i|)^{2^{m+1}}$ when $z_i\ne0,\infty$, the phase in
\eqref{eq:input}.  Also, $h=1$ throughout $[70,82]$, which contains
all the readout supports, independently of the main coordinate.
These are conditional explanations of the formula.  Agreement of the actual and
predicted clocks, and the required bound keeping a register in the unit
disk until its readout, are proved in Section~\ref{sec:orbits} before
the formula is used to analyze the main coordinate.

\begin{figure}[tbp]
\centering
\begin{tikzpicture}[
  phase/.style={draw=black!65,rounded corners=2pt,align=center,
    minimum height=1.25cm,text width=2.8cm,font=\small,inner sep=4pt},
  >={Stealth[length=2mm]}
]
\node[phase,fill=blue!5] (reset) {Reset\\\(t\in[2,5]\)\\\(v'_\ell=0\)};
\node[phase,fill=blue!5,right=0.35cm of reset] (prep)
  {Preparation\\\(t\in(8,12)\)\\add \(U_\ell(t,z)\)};
\node[phase,fill=orange!8,right=0.35cm of prep] (slow)
  {Possible delay\\\(t\in(48,52)\)\\speed set by \(z\)};
\node[phase,fill=blue!5,right=0.35cm of slow] (read)
  {Readout\\\(t\in(74,80)\)\\add aligned terms to \(z^2\)};
\draw[->] (reset)--(prep);
\draw[->] (prep)--(slow);
\draw[->] (slow)--(read);
\draw[|-|,black!85] ([yshift=-0.45cm]prep.south west)
  -- node[below,font=\small,align=center]
  {\(a(t)=L\) throughout preparation, delay and readout}
  ([yshift=-0.45cm]read.south east);
\end{tikzpicture}
\caption{Clock intervals during a complete passage, shown schematically
and not to scale. Values refer to the clock before each update.
All registers reset before preparation. Each channel is read at most
once; for a channel that is read, its preparation signals give a bounded
contribution at that readout.}
\label{fig:clock}
\end{figure}


\subsection{Continuous differentiability}

\begin{proposition}
\label{prop:regularity}
For the choices above, with $q$ finite and $L>\max\{B,2\}$, the series
\eqref{eq:input} defines $C^1$ maps
$U_\ell\colon(\mathbb R/100\mathbb Z)\times\Sigma\to\mathbb C$.
Consequently, \eqref{eq:map} defines a $C^1$ self-map of the nonempty
compact smooth manifold $M$.
\end{proposition}

\begin{proof}
We first estimate the input series in the fixed product Riemannian
metric, then pass to source coordinates.  The chain rule and
$\|DG\|\le B$ give
\begin{equation}
\label{eq:prediction-derivative}
 \|D(G^m)\|\le B^m,\qquad \|DT_m\|\le B^m.
\end{equation}
The second bound uses the real projection $\operatorname{pr}_1$, whose
derivative has norm one everywhere.  In particular, it requires no
choice of output chart for $G^m$.

For an integer $k\ge1$, let
\[
 \Phi_k(z)=\chi(z)(z/|z|)^k\quad(z\in\mathbb C\setminus\{0\}),
 \qquad \Phi_k(0)=\Phi_k(\infty)=0.
\]
These are smooth complex-valued functions.  The support of $\chi$ lies
in a fixed compact annulus, on which the angular map $z\mapsto z/|z|$
and its first derivative are bounded.  The product and chain rules
therefore give a constant $C_\chi$, independent of $k$, such that
\begin{equation}
\label{eq:angular-derivative}
 \|\Phi_k\|_\infty\le1,\qquad
 \|D\Phi_k\|_\infty\le C_\chi(1+k).
\end{equation}
The extensions vanish on fixed neighborhoods of both poles.

Write the $m$th summand, including its prefactor, as
\[
 F_{\ell,m}(t,z)=\frac{L^{1-m}}{10}
             \alpha(t)\beta_\ell(T_m(t,z))\Phi_{2^{m+1}}(z).
\]
Its differential is
\begin{align*}
 DF_{\ell,m}=\frac{L^{1-m}}{10}\bigl(&
 \alpha'(t)\beta_\ell(T_m)\Phi_{2^{m+1}}\,dt\\
 &+\alpha(t)\beta_\ell'(T_m)\Phi_{2^{m+1}}\,DT_m\\
 &+\alpha(t)\beta_\ell(T_m)D\Phi_{2^{m+1}}\bigr).
\end{align*}
Since there are finitely many $\beta_\ell$, their first derivatives have
a common bound.  Equations~\eqref{eq:prediction-derivative}
and~\eqref{eq:angular-derivative} yield a constant $C$, independent of
$\ell$ and $m$, with
\begin{equation}
\label{eq:input-derivative-bound}
 \|F_{\ell,m}\|_\infty+\|DF_{\ell,m}\|_\infty
 \le C L^{1-m}\bigl(1+B^m+2^{m+1}\bigr).
\end{equation}
The prediction and angular derivatives occur in separate product-rule
terms; their bounds are added.  Each of the three resulting geometric
series converges by \eqref{eq:gain-choice}.

For completeness, cover $[0,100]\times\Sigma$ by finitely many relatively
compact source coordinate neighborhoods $V_j$ whose closures lie in larger coordinate
charts $W_j\subset\mathbb R\times\Sigma$.  If $\kappa_j$ is the chart
on $W_j$, the norm of $D\kappa_j^{-1}$ on
$\kappa_j(\overline{V_j})$ is bounded by a constant $C_j$ independent of
$m$.  Thus the coordinate derivatives of
$F_{\ell,m}\circ\kappa_j^{-1}$ satisfy \eqref{eq:input-derivative-bound}
with the additional fixed factor $C_j$.  The series of functions and
of all their first coordinate derivatives consequently converge
uniformly on every such neighborhood.  The usual differentiation
theorem for series now applies: on a smaller coordinate ball, integrate
the first derivatives of the partial sums along line segments and pass
to the uniform limit.  This identifies the derivative of the limit
with the continuous limit of the derivatives.  Hence each $U_\ell$ is
$C^1$ on the fundamental strip.

This proof also covers data with arbitrarily late predicted visits.
At a datum for which $\beta_\ell(T_m)=0$ for every $m$, each summand and
its first derivative vanish: a differentiable nonnegative function has
derivative zero at any zero.  Uniform convergence of the derivative
tails ensures continuity even when nearby nonzero summands have indices
tending to infinity.  No regularity of a selected hitting time is needed.

Because $\alpha$ vanishes near $0$ and $100$, the input and all its
first derivatives vanish on neighborhoods of the periodic seams.
The periodic extension is therefore $C^1$.  Because $\chi$ vanishes
near zero and infinity, the same holds in fixed neighborhoods of the
two poles, so the angular expression introduces no coordinate
singularity.  We also have the uniform bound
\begin{equation}
\label{eq:input-size}
 |U_\ell(t,z)|\le\frac1{10}\sum_{m=1}^{\infty}L^{1-m}
             =\frac{L}{10(L-1)}.
\end{equation}

The clock component of \eqref{eq:map} is smooth and well-defined modulo
$100$.  Each register component is $C^1$ and takes its values in the
fixed disk of radius $3L+L/(10(L-1))$ in the affine chart.  This includes
all source values $v_\ell=\infty$, since $S$ is smooth and zero near
infinity.

It remains to check the main component at $z=\infty$.  Set
\[
 C_0(t,v_1,\ldots,v_q)=20\sum_{\ell=1}^q\bar\beta_\ell(t)S(v_\ell),
 \qquad |C_0|\le60q.
\]
This is a smooth function of all its variables.  In the source inverse
coordinate $w=1/z$ and the corresponding inverse coordinate for the
output, the main component is
\begin{equation}
\label{eq:infinity-chart}
 w'=\frac{w^2}{1+C_0w^2}.
\end{equation}
For $|w|<[2(60q+1)]^{-1/2}$ the denominator has modulus greater than
$1/2$, uniformly in $t$ and the registers.  Formula~\eqref{eq:infinity-chart}
therefore extends smoothly to $w=0$, with value zero.  At all finite
$z$, including $z=0$, the main component is already smooth in the
affine chart.  These checks establish the asserted global $C^1$
regularity of $f$.
\end{proof}
